% chapters2/ch-rmp.tex — the review-article test: figures of the
% matrix-product-state review arXiv:2011.12127 reproduced as
% formula-first worked examples.  Every body below is adapted from the
% verified reproduction corpus (hard01..hard12 v3, plane acceptance);
% manual2.tex inputs this file directly.
\section{The review-article test}\label{ch:rmp}

A notation earns its keep by reproducing the figures of a paper it
was not designed around.  This chapter works through diagrams of the
matrix-product-state review of Cirac, P\'erez-Garc\'ia, Schuch and
Verstraete,\sidenote{Rev.\ Mod.\ Phys.\ \textbf{93}, 045003 (2021);
arXiv:2011.12127.  The examples follow the review's Sec.~II--III:
canonical forms, fixed points, blocking, matrix product operators, the
tangent space, and PEPS.} chosen because each one leans on a different
part of the language: cups and canonical triangles, on-bend fixed-point
matrices, per-column legs on wide gates, tall gauges spanning two
layers, per-site physical traces, and the oblique two-dimensional
sheets.  Wherever an equation is drawn, the margin quotes the
boundary signature that both sides must share.

\subsection{Canonical forms}

A uniform MPS tensor $A_L$ is \emph{left-canonical} when its matrices
are jointly isometric from the left.  Contracting the shared left
virtual index and the physical index of a ket--bra pair must give the
identity on the right pair --- so both sides of the drawn equation keep
exactly two east stubs, and nothing else, open.

\begin{tnexample}[
    formula   = {\sum_s (A_L^s)^\dagger A_L^s = \mathbb{I}},
    caption   = {left orthonormality of $A_L$},
    index     = {The triangle is the left-canonical tensor; its apex
                 points toward the surviving index.  The vertical bond
                 is the summed physical index $s$; the west cup is the
                 summed left virtual index $\alpha$.  Both sides keep
                 the pair $(\beta,\beta')$ open: the equation is
                 between matrices.},
    signature = {boundary|virtual-west=0|virtual-east=2|physical-up=0|physical-down=0},
    label     = {ex:rmp-ortho-left},
  ]
% LHS: ket A_L over bra conj(A_L); sandwich pairs the physical
% index s; west=cup contracts the shared left
% virtual index alpha
\[
  \begin{tenkz}[sandwich, west=cup]
    \tn[tri=l]{A_L}\\
    \tn*[tri=l]{A_L}
  \end{tenkz}
  \;=\;
  % RHS: the identity on the open (beta, beta') pair -- a bare
  % cup; the ghosts are the wire endpoints
  \begin{tenkz}[rows={wire,wire}, west=cup]
    \tnghost{}\\
    \tnghost{}
  \end{tenkz}
\]
\end{tnexample}

The right-canonical condition is the mirror image, and so is its
source: the triangles flip to \tnval{tri=r}, the cup moves to the east
side, and the identity keeps the west pair open.

\begin{tnexample}[
    formula   = {\sum_s A_R^s (A_R^s)^\dagger = \mathbb{I}},
    caption   = {right orthonormality of $A_R$},
    index     = {Roles of the sides swap: the east cup is now the
                 summed right index $\beta$, and the open pair
                 $(\alpha,\alpha')$ sits west on both sides.  The
                 review draws both equations in isometry orientation,
                 apex up, contraction arcs below --- a mirror image
                 of this chain geometry, same topology.},
    signature = {boundary|virtual-west=2|virtual-east=0|physical-up=0|physical-down=0},
    label     = {ex:rmp-ortho-right},
  ]
% same sandwich with right-canonical triangles; the cup is the
% contraction over the shared right index beta
\[
  \begin{tenkz}[sandwich, east=cup]
    \tn[tri=r]{A_R}\\
    \tn*[tri=r]{A_R}
  \end{tenkz}
  \;=\;
  % RHS: the identity on the open (alpha, alpha') pair
  \begin{tenkz}[rows={wire,wire}, east=cup]
    \tnghost{}\\
    \tnghost{}
  \end{tenkz}
\]
\end{tnexample}

\subsection{The reduced density matrix}

Write $A^{(s)} = A^{s_1}\!\cdots A^{s_n}$ for an $n$-site word, and let
$\rho^L$, $\rho^R$ be the dominant left and right fixed points of the
transfer map.  The review's $n$-site reduced density matrix is the
racetrack: a ket rail with all up legs open, the conjugate rail below
with all down legs open, and the two rails closed into each other at
both ends by cups carrying the fixed points.

\begin{tnexample}[
    formula   = {[\rho_n]_{(s)(s')} =
                 \mathrm{Tr}\bigl[\rho^L A^{(s)}\rho^R
                 A^{(s')\dagger}\bigr]},
    caption   = {the $n$-site reduced density matrix},
    index     = {\tnval{nopair} keeps every physical leg open ---
                 $\rho_n$ is a matrix on the $n$-site physical space,
                 $n=4$ drawn.  The bends are the closure of the two
                 rails into each other, and the dots riding them are
                 the fixed points $\rho^L$ (west) and $\rho^R$
                 (east); the formula reads clockwise around the
                 racetrack from the west bend.},
    signature = {boundary|virtual-west=0|virtual-east=0|physical-up=4|physical-down=4},
    label     = {ex:rmp-rdm},
  ]
% ket rail: A^{s_1}..A^{s_4}, up legs open; bra rail: the
% conjugates, down legs open; the cups + on-bend dots are the
% trace through rho^L (west, the left fixed point) and rho^R
% (east, the right fixed point)
\[
  \rho_n \;=\;
  \begin{tenkz}[rows={ket:nopair, bra},
                west={cup=$\rho^L$}, east={cup=$\rho^R$}]
    \tn{} & \tn{} & \tn{} & \tn{}\\
    \tn*{} & \tn*{} & \tn*{} & \tn*{}
  \end{tenkz}
\]
\end{tnexample}

The review then asks for the spectrum of this object.  The picture
enters the formula like any other math atom: its wire axis sits on the
math axis, so stretchy delimiters size and centre on the racetrack.

\begin{tnexample}[
    formula   = {\mathrm{eig}(\rho_n)},
    caption   = {the spectrum of $\rho_n$},
    index     = {Same racetrack as Example~\ref{ex:rmp-rdm}; the
                 \tnval{\char`\\left( ... \char`\\right)} pair centres
                 on the picture because the wire axis is the math
                 axis.  eig is algebra around one picture, so there is
                 no side-to-side signature balance to check.},
    signature = {boundary|virtual-west=0|virtual-east=0|physical-up=4|physical-down=4},
    label     = {ex:rmp-eig},
  ]
% rho_n racetrack inside stretchy parentheses: the picture's
% wire axis rides the math axis
\[
  \mathrm{eig}\left(
  \begin{tenkz}[rows={ket:nopair, bra},
                west={cup=$\rho^L$}, east={cup=$\rho^R$}]
    \tn{} & \tn{} & \tn{} & \tn{}\\
    \tn*{} & \tn*{} & \tn*{} & \tn*{}
  \end{tenkz}
  \right)
\]
\end{tnexample}

\subsection{Blocking to a fixed point}

Under renormalization a fixed-point MPS reproduces itself when two
sites are blocked into one: the blocked pair equals a single tensor
followed by an isometry $U$ on the doubled physical index.  Both sides
must show two up legs --- one per blocked site --- which is exactly
what \tnkey{legs at=} provides on a wide glyph.

\begin{tnexample}[
    formula   = {A^{i_1}A^{i_2} = \sum_j U_{(i_1 i_2),j}\,A^j},
    caption   = {the blocking fixed-point equation},
    index     = {$U$ is a pill of width two; \tnval{legs
                 at=\{1,2\}} opens one up leg over each spanned
                 column, so both sides read physical-up $=2$.  The
                 vertical $U$--$A$ bond is the summed index $j$;
                 \tnval{op:none} seals $U$'s virtual boundary ---
                 the isometry has no horizontal wires.},
    signature = {boundary|virtual-west=1|virtual-east=1|physical-up=2|physical-down=0},
    label     = {ex:rmp-blocking},
  ]
% LHS: A^{i_1} A^{i_2} -- two MPS tensors, one open up leg each
\[
  \begin{tenkz}[physical=up, tensor style=box]
    \tn{A} & \tn{A}
  \end{tenkz}
  \;=\;
  % RHS: U keeps one leg per blocked site (legs at={1,2}); the
  % vertical U-A bond is the summed physical index j
  \begin{tenkz}[rows={op:none, ket}, tensor style=box]
    \tn[pill, wide=2, legs at={1,2}]{U} & \\
    \tn[wide=2]{A} &
  \end{tenkz}
\]
\end{tnexample}

\subsection{Purification and zero correlation length}

A positive matrix product operator admits a local purification: the
MPO tensor $O$ factors, up to a gauge $X$ on the virtual index, into a
purification tensor $A$ glued to its own conjugate.  The review draws
$X$ as a tall capsule spanning both layers with the label inscribed,
which is \tncmd{tnX} with \tnkey{wires=2}.

\begin{tnexample}[
    formula   = {O = X^{-1}\bigl[{\textstyle\sum_k}
                 A^k\!\otimes\overline{A}^k\bigr]X},
    caption   = {local purification of a positive MPO},
    index     = {The $A$--$\overline{A}$ bond is the purification
                 index $k$ (the automatic pairing of the two
                 \tnval{op} rows); $s$ (up) and $s'$ (down) stay
                 open.  Each capsule is one atom spanning both layer
                 wires; \tnval{op:none} seals the top row, so exactly
                 one virtual index leaves each capsule --- the
                 signature matches $O$'s.  One recorded gap: the
                 review's outer wire enters the capsule at its
                 vertical centre (the fused index); the grid opens
                 one row's boundary, so here it sits at the lower
                 row's height.},
    signature = {boundary|virtual-west=1|virtual-east=1|physical-up=1|physical-down=1},
    label     = {ex:rmp-purification},
  ]
% LHS: the 4-leg MPO tensor O (s up, s' down, virtual a/b)
\[
  \begin{tenkz}[physical=updown, tensor style=box]
    \tn{O}
  \end{tenkz}
  \;=\;
  % RHS: X^{-1} [ A (x) conj(A) ] X -- capsules span both
  % layers; the A-conj(A) bond is the purification index k;
  % align=2 puts the open row on the math axis
  \begin{tenkz}[rows={op:none, op}, align=2, tensor style=box]
    \tnX[wires=2]{X^{-1}} & \tn{A} & \tnX[wires=2]{X}\\
                          & \tn*{A} &
  \end{tenkz}
\]
\end{tnexample}

At the opposite extreme sit the zero-correlation-length density
operators.  With $E = \sum_s M^{ss}$ --- the MPDO tensor with its own
physical pair closed by a per-site trace --- the condition is that $E$
be idempotent.  \tnkey{trace=physical} draws that per-site closure.

\begin{tnexample}[
    formula   = {E\,E = E,\quad E = {\textstyle\sum_s} M^{ss}},
    caption   = {zero correlation length for an MPDO},
    index     = {\tnval{trace=physical} arcs each site's up leg
                 around the glyph into its own down leg: the per-site
                 physical trace.  Traced legs are closed indices and
                 leave the signature, so both sides read
                 $(1,1,0,0)$.},
    signature = {boundary|virtual-west=1|virtual-east=1|physical-up=0|physical-down=0},
    label     = {ex:rmp-zcl},
  ]
% LHS: E.E -- two physically traced tensors, one virtual bond
\[
  \begin{tenkz}[physical=updown, trace=physical,
                tensor style=box]
    \tn[mpo]{M} & \tn[mpo]{M}
  \end{tenkz}
  \;=\;
  % RHS: E -- one physically traced tensor
  \begin{tenkz}[physical=updown, trace=physical,
                tensor style=box]
    \tn[mpo]{M}
  \end{tenkz}
\]
\end{tnexample}

\subsection{The tangent space}

In the gauge where the chain is left-canonical up to site $i$ and
right-canonical after it, the projector onto the tangent space of a
uniform MPS is
\[
  P_{T(A)} \;=\; \sum_{i=1}^{n}
    P_L^{(i-1)}\otimes\mathbb{I}_i\otimes P_R^{(i+1)}
  \;-\; \sum_{i=1}^{n} P_L^{(i)}\otimes P_R^{(i+1)},
\]
where $P_L^{(k)}$ projects onto the states reached by an $A_L$-chain
on the sites up to $k$, and $P_R^{(k)}$ mirrors it from the right.
Each block projector is an all-legs-open window of canonical triangles
whose edge cup is the bracket
$\sum_\alpha|\Phi_\alpha\rangle\langle\Phi_\alpha|$; pictures set side
by side are tensor factors, so a summand is assembled exactly as the
formula factors it.

\begin{tnexample}[
    formula   = {P_L^{(i-1)}\otimes\mathbb{I}_i\otimes
                 P_R^{(i+1)}},
    caption   = {one summand of the tangent-space projector},
    index     = {Left block: an $A_L$ window, all physical legs open
                 (\tnval{nopair}), its east cup the bracket at the
                 block edge; the right block mirrors it.  The
                 site-$i$ identity is algebra, not ink: a cup cannot
                 sit at a hole's edge, so the summand carries
                 $\mathbb{I}_i$ symbolically.},
    signature = {boundary|virtual-west=2|virtual-east=0|physical-up=2|physical-down=2},
    label     = {ex:rmp-tangent},
  ]
% P_L^(i-1): all-legs-open A_L window; the east cup is the
% tangent-space bracket sum_a |Phi_a><Phi_a|
\[
  \begin{tenkz}[rows={ket:nopair, bra}, east=cup]
    \tndots & \tn[tri=l]{} & \tn[tri=l]{}\\
    \tndots & \tn*[tri=l]{} & \tn*[tri=l]{}
  \end{tenkz}
  % 1_i: algebra, not ink -- no bare-wire atom can be bracketed
  \otimes \mathbb{I}_i \otimes
  % P_R^(i+1): the mirrored A_R window, bracket at the west
  \begin{tenkz}[rows={ket:nopair, bra}, west=cup]
    \tn[tri=r]{} & \tn[tri=r]{} & \tndots\\
    \tn*[tri=r]{} & \tn*[tri=r]{} & \tndots
  \end{tenkz}
\]
\end{tnexample}

The algebraic $\mathbb{I}_i$ contributes no drawn legs, so the two
sums balance as drawn: each summand's ink totals
$(2,2,4,4)$,\sidenote{Per picture: the line quoted beside
Example~\ref{ex:rmp-tangent} is a left block's; a right block mirrors
it, \tnlogline{virtual-east=2} in place of
\tnlogline{virtual-west=2}.} whether or not the brackets meet back to
back --- the second summand is the same two blocks with $\otimes$
alone between them.

\subsection{Two dimensions}

A PEPS assigns one five-index tensor to each site of a square lattice:
four virtual bonds to the lattice neighbours and one physical leg.
The review draws finite windows obliquely so the physical legs leave
the lattice plane visibly; \tnval{frame=oblique} shears the bonds
only, keeping dots, labels, and legs in the paper frame.

\begin{tnexample}[
    formula   = {|\Psi\rangle = \sum_{\mathbf{s}}
                 \mathcal{C}\bigl(\{A^{s_{xy}}\}\bigr)\,
                 |\mathbf{s}\rangle},
    caption   = {a finite PEPS window, drawn obliquely},
    index     = {Each dot is a site tensor $A^{s_{xy}}$; the in-sheet
                 bonds are its four virtual indices, and
                 \tnval{boundary legs} lets the bonds cut by the
                 window dangle.  The short vertical stub at each site
                 is the open physical index, readable against the
                 sheared lattice.},
    label     = {ex:rmp-peps},
  ]
% a 4x4 PEPS window: sheared bonds, per-site physical legs,
% cut bonds dangling at the window boundary; the marked central segment is
% one representative virtual contraction and changes no index ownership
\begin{tenkzlattice}[rows=4, cols=4, frame=oblique,
                     physical=up, boundary legs]
  \tnedge[role=marked]{(2,2)-(2,3)}
\end{tenkzlattice}
\end{tnexample}

Expectation values double the sheet: a ket layer over its conjugate,
with a site-wise pairing leg contracting each physical index.  Leaving
the pairing legs of chosen cells dangling turns the double layer into
a marginal on those cells.

\begin{tnexample}[
    formula   = {\rho_C = \operatorname{Tr}_{\overline{C}}\,
                 |\Psi\rangle\langle\Psi|},
    caption   = {the double layer with open columns},
    index     = {Two sheets, ket over bra; each vertical pairing leg
                 is one contraction $\sum_s A^s\overline{A}^s$.
                 Cells listed in \tnval{open=} keep their pairing
                 legs dangling --- here columns 2 and 4, so the
                 picture is the marginal $\rho_C$ of those columns.},
    label     = {ex:rmp-doublelayer},
  ]
% ket-over-bra double layer; columns 2 and 4 stay open, so the
% object is the marginal of those columns
\begin{tenkzplanes}[rows=4, cols=4, sheets={ket, bra},
                    open={(1-4,2),(1-4,4)}]
\end{tenkzplanes}
\end{tnexample}

Ten review figures, and the load was carried by a handful of keys:
cups with on-bend matrices, \tnval{nopair}, \tnkey{legs at=},
\tnkey{wires=2}, \tnval{trace=physical}, \tnval{frame=oblique}, and
the double-layer environment.  What the test could not draw --- the
bare identity wire inside a bracketed summand, the centred fused wire
on a tall gauge --- is recorded here beside the figures that need it,
and again in the roadmap of \S\ref{ch:trouble}.
